% --------------------------------- Preamble -----------------------------------
\documentclass[11pt, letterpaper, twocolumn]{article}

\usepackage[T1]{fontenc} % use 8-bit fonts
\usepackage[margin=0.75in, columnsep=0.2in, % 3.4 in columns
    headsep=0.28in, footskip=0.42in]{geometry} % paper sizes
\usepackage{amsmath, amsfonts, amssymb} % improved math
\usepackage{bm} % bold math (must be after amsmath)
\usepackage{upquote} % upright single quotes in verbatim
\usepackage{pgfplots} % includes tikz, xcolor, graphicx
\usepackage[colorlinks=true, allcolors=gray]{hyperref}

\raggedbottom % do not force stretching text to fill the page
\setlength{\emergencystretch}{\hsize} % avoid overfull lines
\overfullrule=1mm % black band to show overfull lines
\pdfsuppresswarningpagegroup=1 % remove pdf image warnings
\frenchspacing % consistent spacing between words
\newcommand{\ul}[1] % vectors
    {{}\mkern1mu\underline{\mkern-1mu#1\mkern-1mu}\mkern1mu}
\DeclareSymbolFont{euler}{U}{eur}{m}{n} % symbol font
\DeclareMathSymbol{\PI}{\mathord}{euler}{25} % up pi
\pgfplotsset{compat=newest} % use newest pgfplot version
\newcommand{\EE}[2]{\ensuremath{{{#1}\!\times\!10^{#2}}}}
% ------------------------------------------------------------------------------

\title{Particle Filter Localization with Lidar}
\author{}
\date{}

\begin{document}
\maketitle

\section*{Description}

This project implements a particle filter for robot localization in a mapped
indoor environment. The robot state consists of planar position and heading, 
and the filter uses control inputs together with noisy lidar range measurements
to estimate the robot pose over time. The particle filter represents the state
distribution using weighted samples, making it well suited for nonlinear
systems and non-Gaussian uncertainty.

\subsection*{Problem Formulation}

The robot state is defined as
\[
x = \begin{bmatrix} p_x \\ p_y \\ \theta \end{bmatrix}
\]
where $p_x$ and $p_y$ are the 2D position coordinates and $\theta$ is the heading
angle.

The input consists of the commanded speed $s$ and heading rate $\omega$, while 
the measurement is a scalar lidar range to the nearest wall. Ground-truth states,
inputs, and measurements are provided from the recorded data. The goal is to 
estimate the robot pose over time using a particle filter.

\subsection*{System Model}

The continuous-time dynamics are given by 
\[\begin{aligned}
\dot{p}_x &= (s+\nu_s)\cos(\theta) \\
\dot{p}_y &= (s+\nu_s)\sin(\theta) \\
\dot{\theta} &= \omega + \nu_{\omega}
\end{aligned}
\]
where $\nu_s$ and $\nu_{\omega}$ are process noise terms.

Using forward Euler integration with sampling period $T$, the discrete-time model becomes
\[\begin{aligned}
p_{x,k+1} &= p_{x,k} + (s_k + \nu{s,k})\cos(\theta_k)T \\
p_{y,k+1} &= p_{y,k} + (s_k + \nu{s,k})\sin(\theta_k)T \\
\theta_{k+1} &= \theta_k + (\omega_k + \nu_{\omega, k})T
\end{aligned}
\]
The lidar measurement model is nonlinear and is evaluated using the provided map-based
function $h(x,k)$.

\subsection*{Initial Conditions}

We define the following covariance matrices
\[\begin{aligned}
P_0 &= \begin{bmatrix}
    {\sigma_x}^2 & 0 & 0 \\
    0 & {\sigma_y}^2 & 0 \\
    0 & 0 & {\sigma_z}^2
\end{bmatrix} &\quad Q &= \begin{bmatrix}
    {\sigma_s}^2 & 0 \\
    0 & {\sigma_{\omega}}^2
\end{bmatrix} \\ 
R &= {\sigma_r}^2 &\quad \sigma_r &= 10 cm \\
\sigma_x &= 0.1 cm &\quad \sigma_s &= 32 cm/s \\
\sigma_y &= 0.1 cm &\quad \sigma_{\omega} &= 0.55 rad/s \\
\sigma_a &= 5\times10^{-5} rad   
\end{aligned}  
\]
Where $P_0$ is the initial state covariance matrix, $Q$ is the process noise covariance
matrix, and $R$ is the measurement noise. Each $\sigma^2$ is the covariance term for 
its respective state/input/measurement.

\section*{Particle Filter Method}

The state distribution is represented by a set of particles and associated weights. The
particles are initalized from a Gaussian distribution with mean equal to the initial 
true state and covariance determined by the assumed uncertainty in position and heading.

At each time step, the filter performs the following steps:

\textbf{Update:} The particle weights are updated using the likelihood of the lidar 
measurement under a Gaussian measurement noise model.

\textbf{Estimate:} The state estimate is computed as the weighted average of the particles,
and the state covariance is estimated from the weighted particle deviations.

\textbf{Resample:} When the effective number of particles becomes too small, resampling is
performed to eliminate low-weight particles and duplicate high-weight particles.

\textbf{Propagate:} The particles are advanced through the nonlinear dynamics using the 
measured inputs and independently sampled process noise.

\section*{Simulation}

The particle filter is evaluated under four different initialization conditions:

\begin{enumerate}
    \item Known starting position and known heading (10 particles)
    \item Unknown starting position and known heading (25 particles)
    \item Known starting position and unknown heading (40 particles)
    \item Unknown starting position and unknown heading (650 particles)
\end{enumerate}

When the starting position is known, the initial position standard deviations are set
to nominal values. When the starting position is unknown, the initial standard deviations
in $x$ and $y$ are increased to 100 cm. When the starting heading is unknown, the 
initial heading standard deviation is set to $\pi$ rad.

These cases illustrate how the required number of particles increases as the initial
uncertainty becomes larger.

\pagebreak

\section*{Results}

\subsection*{Path Plots}

Figures~\ref{fig:paths1}--\ref{fig:paths4} show the true and estimated robot trajectories
for the four initialization cases. \\
    
\begin{figure}[h]
    \centering
    \includegraphics[width=\linewidth]{../figures/fig_paths_1.pdf}
    \caption{Known starting position and known heading.}
    \label{fig:paths1}
\end{figure}

\begin{figure}[h]
    \centering
    \includegraphics[width=\linewidth]{../figures/fig_paths_2.pdf}
    \caption{Unknown starting position and known heading.}
    \label{fig:paths2}
\end{figure}

\pagebreak

\begin{figure}[h]
    \centering
    \includegraphics[width=\linewidth]{../figures/fig_paths_3.pdf}
    \caption{Known starting position and unknown heading.}
    \label{fig:paths3}
\end{figure}

\begin{figure}[h]
    \centering
    \includegraphics[width=\linewidth]{../figures/fig_paths_4.pdf}
    \caption{Unknown starting position and unknown heading.}
    \label{fig:paths4}
\end{figure}

\pagebreak

\subsection*{Error Plots}

For the cold-start case, Figures~\ref{fig:xerr}--\ref{fig:aerr} show the estimation 
errors for position and heading.

\begin{figure}[h]
    \centering
    \includegraphics[width=\linewidth]{../figures/fig_x_error.pdf}
    \caption{$x$-position estimation error.}
    \label{fig:xerr}
\end{figure}

\begin{figure}[h]
    \centering
    \includegraphics[width=\linewidth]{../figures/fig_y_error.pdf}
    \caption{$y$-position estimation error.}
    \label{fig:yerr}
\end{figure}

\begin{figure}[h]
    \centering
    \includegraphics[width=\linewidth]{../figures/fig_a_error.pdf}
    \caption{Heading estimation error.}
    \label{fig:aerr}
\end{figure}

\pagebreak

\subsection*{Observations}

Starting with the two cases that are the most intuitive, the first case, 
with a known starting position and known heading performed very well, even
with just 10 particles. This is because it is quite easy to track on object
with low noise, even with only 10 particles. 

The second intuitive case is the fourth case, with both an unknown starting
position and an unknown heading. The filter performed well overall, but took
the longest to converge on the truth. With a large region of position uncertainty
and a large range of heading possibilities, it takes a large number of
particles to encapsulate the truth. However, after the measurements come in,
the particles closest to truth will begin to be weighted much higher and
the filter will quickly converge (assuming good measurements).

The cases that are fairly unintuitive are cases 2 and 3, where only one
of the position and heading is unknown. Case 2, with an unknown position
and a known heading, consistently produced the highest error. Taking a longer
amount of time to converge, on average, than case 4 seems counterintuitive,
but diving deeper makes sense. The initial uncertainty of the robots
position is just as large as case 4, but with far fewer particles. As a 
result, the particles are spread too thin to find the robot's true location.
The filter then latches onto incorrect position estimates, which in turn
corrupt the "known" heading estimate as the filter tries to explain
measurement errors. 

The last case is case 3, with a known position and unknown heading. I orginally
thought this would be a higher error case, but it turned out to be the lowest
on average. I underestimated the power of adding more particles, and 40
turned out to be more than enough to cover the range of possible headings.
With the position known, the lidar measurements provide very clear information,
allowing the filter to quickly and accurately determine the correct heading.
However, as you decrease the number of particles, this case causes
a divergence of the filter quite often.

Most of what we discussed above was about the initial conditions of the
particle filter and how long each took to converge. One important thing
to note is that after convergence, the filters with more particles estimated
the true position of the robot more accurately. This seems fairly self-explanatory.

On the topic of resampling, it is also important to note the effective
sample values I found to be most accurate. I ended up keep Je = 0.6, so 
if the effective number of samples drops below 0.6, I would resample. This
number seems reasonable in that it is not too high that we let the particles
spread out too much, and not too low that we don't let them spread out enough.
On average, this number looked to provide the highest accuracy in each case.
Lowering or increasing the number significantly caused the filter to 
diverge frequently, as expected.

\section*{Discussion}

The results show that particle filters can successfully localize a robot in a 
nonlinear environment using noisy control inputs and lidar measurements. When
the initial uncertainty is small, relatively few particles are needed because
the filter only needs to represent a narrow region of the state space. As the
initial uncertainty increases, many more particles are required to adequately
cover the possible robot poses.

From the animation, the particles tended to spread out in regions where the 
lidar measurements were less informative or where multiple poses produced
similar expected measurements. The particles coalesced when the robot encountered
geometrically informative parts of the environment, allowing the measurement
udpate to strongly favor a smaller subset of poses.

This behavior highlights one of the main strengths of the particle filter: it
can represent multi-hypothesis uncertainty and then concentrate around the 
correct solution as informative measurements become available.

\end{document}